% vim: set textwidth=78 autoindent:
% options for boxes provided by fancybox package are:
% \fbox, \shadowbox, \doublebox, \ovalbox and \Ovalbox
%%%%%%%%%%%%%%%%8%%%%%%%%%%%%%%%%%%%%%%%%%%%%%%%%%%%
%
% general styles for text that the user should type in from the keyboard
\renewcommand{\usertext}[1]{\texttt{#1}}
\renewcommand{\filename}[1]{\texttt{#1}}
% style for user input by individual keystrokes
\renewcommand{\keystroke}[1]{\cornersize{.6}\ovalbox{\textsf{#1}}}
% generic style for text that appears in the GUI
\newcommand{\guilabel}[1]{\textsf{#1}}
\newcommand{\guiheading}[1]{\textsf{#1}}
% some styles for use only when discussing Python coding
\renewcommand{\classname}[1]{\textsf{\textbf{#1}}}
\renewcommand{\method}[1]{\textsf{\textit{#1}}}
%\renewcommand{\fieldname}[1]{\textsl{#1}}
\newcommand{\object}[1]{\textsf{\textit{#1}}}
% some styles for the Qt GUI
\newcommand{\qtmainmenuopt}[1]{\textsf{#1}}
\newcommand{\qtdropmenuopt}[1]{\textsf{#1}}
\newcommand{\qtdialog}[1]{\textsf{#1}}
% ??? these styles are from Gary
\renewcommand{\server}[1]{\textit{#1}}
\renewcommand{\sqltable}[1]{\textsf{\textbf{#1}}}


%
\newcommand{\dialog}[1]{
\fcolorbox[rgb]{0.2,0.2,0.95}[rgb]{0.2,0.2,0.95}{%
\textbf{\textcolor{white}{#1}}}}
%
\setlength{\shadowsize}{2pt}%
% nix, win and osx for platform-specific documentation
% old versions
%\newcommand{\nix}[1]{Under GNU/Linix, #1}
%\newcommand{\win}[1]{Under MS Windows, #1}
%\newcommand{\osx}[1]{Under Mac OS X, #1}

% add operation system icons to figure \caption
% usage \caption{Text \wincaption}
\newcommand{\nixcaption}{\protect\includegraphics[height=4mm]{nix.eps}}
\newcommand{\wincaption}{\protect\includegraphics[height=4mm]{win.eps}}
\newcommand{\osxcaption}{\protect\includegraphics[height=4mm]{osx.eps}}

% new versions
\newcommand{\nix}[1]{\includegraphics[height=5mm]{nix.eps} #1}
\newcommand{\win}[1]{\includegraphics[height=5mm]{win.eps} #1}
\newcommand{\osx}[1]{\includegraphics[height=5mm]{osx.eps} #1}

% button
\renewcommand{\button}[1]{%
\raisebox{-6pt}{%
\shadowbox{\guilabel{#1}}
}}
% mainmenuopt
\newcommand{\mainmenuopt}[1]{%
\raisebox{-6pt}{%
\setlength{\fboxsep}{0pt}%
\shadowbox{\setlength{\fboxsep}{2pt}%
\fcolorbox[gray]{0.9}[gray]{0.9}%
{\guilabel{#1}}%
}}}
% dropmenuopt
\newcommand{\dropmenuopt}[1]{%
\raisebox{-6pt}{{%
\setlength{\fboxsep}{0pt}%
\shadowbox{\setlength{\fboxsep}{2pt}%
\fcolorbox[rgb]{.95,.95,0.8}[rgb]{.95,.95,0.8}%
{ \guilabel{#1}}}%
}}}
% dropmenucheck
\newcommand{\dropmenucheck}[1]{%
\raisebox{-6pt}{{%
\setlength{\fboxsep}{0pt}%
\shadowbox{\setlength{\fboxsep}{2pt}%
\fcolorbox[rgb]{.95,.95,0.8}[rgb]{.95,.95,0.8}%
{ $\boxtimes$ \guilabel{#1}}}%
}}}
% dropmenuopttwo
\newcommand{\dropmenuopttwo}[2]{%
\raisebox{-6pt}{%
{%
\setlength{\fboxsep}{0pt}%
\shadowbox{\setlength{\fboxsep}{2pt}%
\fcolorbox[rgb]{.95,.95,0.8}[rgb]{.95,.95,0.8}%
{\includegraphics[width=3mm]{#1} \guilabel{#2}}}%
}}}
% tooltip
\newcommand{\tooltip}[1]{%
\raisebox{-2pt}{%
\fcolorbox{black}[rgb]{1,1,0.8}{\guilabel{#1}}%
}}
% toolbtntwo
\newcommand{\toolbtntwo}[2]{%
\raisebox{-6pt}{%
\shadowbox{\includegraphics[width=7mm]{#1}}} %
\tooltip{#2}%
}
% toolboxtwo
\newcommand{\toolboxtwo}[2]{%
\raisebox{-6pt}{%
\shadowbox{%
\raisebox{-2pt}{\includegraphics[width=5mm]{#1}}%
\guilabel{ #2}%
}}}
% toolboxthree
\newcommand{\toolboxthree}[3]{%
\raisebox{-6pt}{%
\shadowbox{%
\includegraphics[width=5mm]{#1}%
\includegraphics[width=5mm]{#2}%
 \guilabel{#3}%
}}}
% toolboxfour
\newcommand{\toolboxfour}[4]{%
\raisebox{-6pt}{%
\shadowbox{%
\includegraphics[width=5mm]{#1}%
\includegraphics[width=5mm]{#2}%
\includegraphics[width=5mm]{#3}%
 \guilabel{#4}%
}}}
%not to be used anymore
%\newcommand{\toolbox}[2]{%
%\toolboxtwo{#1}{#2}%
%}
% tab
\newcommand{\tab}[1]{%
\raisebox{-6pt}{%
\shadowbox{\guilabel{#1}}
}}
% checkbox
\newcommand{\checkbox}[1]{%
\raisebox{2pt}{%
\fbox{%
\raisebox{-4pt}{%
\shadowbox{x}} %
\guilabel{#1}}%
}}
% radiobuttonoff
\newcommand{\radiobuttonoff}[1]{%
\raisebox{-4pt}{%
\setlength{\fboxsep}{1pt}%
\shadowbox{%
$\bigcirc$%
}} %
\guilabel{#1}%
}
% radiobuttonon
\newcommand{\radiobuttonon}[1]{%
$\odot$ %
\guilabel{#1}%
}
% selectnumber
\newcommand{\selectnumber}[2]{%
\fbox{%
{#1} \fbox{{#2} %
\raisebox{-6pt}{%
\setlength{\fboxsep}{1pt}%
\shadowbox{%
${\blacktriangle}\atop{\blacktriangledown}$%
}}}}}
%
% selectstring
\newcommand{\selectstring}[2]{%
\fbox{%
\guilabel{#1} \fbox{\guilabel{#2} %
\raisebox{-2pt}{%
\setlength{\fboxsep}{1pt}%
\shadowbox{%
$\blacktriangledown$%
}}}}}
%
% browsebutton
\newcommand{\browsebutton}{%
\raisebox{-6pt}{%
\shadowbox{\rule[-1mm]{0mm}{4mm}{$\ldots$}}
}}
%
% selectcolor
\newcommand{\selectcolor}[2]{%
\fbox{\guilabel{#1} %
\raisebox{-6pt}{%
\setlength{\fboxsep}{0pt}%
\shadowbox{\setlength{\fboxsep}{2pt}%
\fcolorbox{#2}{#2}{\rule{0mm}{5mm}\rule{35mm}{0mm}%
}}}}}
%
% slider
\newcommand{\slider}[3]{%
\fbox{%
\guilabel{#1} \guilabel{#2}\% %
\raisebox{-2pt}{%
\setlength{\fboxsep}{1pt}%
\shadowbox{%
$\triangledown$%
}}%
\negthinspace\rule[1mm]{20mm}{1mm}
}}
%
% input text
\newcommand{\inputtext}[2]{%
\fbox{%
\guilabel{#1} %
\raisebox{-6pt}{%
\shadowbox{\usertext{#2}}%
}}}
%
%
\subsection{Conventions}\label{label_conventions}

The conventions used in this manual are as follows. 

\minisec{GUI Conventions}
\begin{itemize}
%
%Use \mainmenuopt for main menu items that have no icon only text
%Main Menu includes: File, View, Layer, Settings, Plugins, Help
\item Menu Options: \mainmenuopt{Layer} > %
%
%Use \dropmenuopttwo for a drop-down menu item with an icon
%Use \dropmenuoptone for a drop-down menu item with no icon 
\dropmenuopttwo{mActionAddRasterLayer}{Add a Raster Layer}

or

\mainmenuopt{View} > %
\dropmenuopt{Toolbar Visibility} > \dropmenucheck{Digitizing}
%
%Use \toolbtntwo for the toolbar items, including those that open dialogs
%These have an icon, and display a tooltip on hover
%Its really important to get the icon in there because that's what a user
%has to search for. The tooltip adds confirmation
\item Tool: \toolbtntwo{mActionAddRasterLayer}{Add a Raster Layer}
%
%Use \button for a clickable button which has no icon, just text
%Save As Default is a button that appears in the Layer Properties dialog.
\item Button: \button{Save as Default}
%
%
%Use \dialog for the title of a dialog box
%such as the Layer Properties dialog.
%Please use the exact title as it appears in the GUI.
\item Dialog Box Title: \dialog{Layer Properties}
%
%Use \tab for clickable tabs which have no icons, just text
%General is a tab that appears in the Layer Properties dialog.
%At the moment, it looks just like \button, but that may change.
\item Tab: \tab{General}
%
%Use \toolboxtwo, \toolboxthree or \toolboxfour
% for GRASS toolbox (not toolbar) items
%These are what you see in the menu after you click Open GRASS Tools
%The one you need depends on how many icons are required
%The number part of the name is the total number of arguments (N)
%The number of icons is (N-1)
\item Toolbox Item: \toolboxtwo{nviz.1.eps}{nviz - Open 3D-View in NVIZ}
%
%Use \checkbox for a checkbox item in a dialog popup
\item Checkbox: \checkbox{Render}
%
%Use \radiobuttonoff for a radio button item in a dialog popup
\item Radio Button:  \radiobuttonon{Postgis SRID} \radiobuttonoff{EPSG ID}
%
% Use \selectnumber for a selection box with up and down arrows
% and a numerical value
\item Select a Number: \selectnumber{Hue}{60}
%
% Use \selectstring for a selection box with down arrows
% and a string value
\item Select a String: \selectstring{Outline style}{---Solid Line}
%
%
% Use \browsebutton for a button that opens a file browser popup
\item Browse for a File: \browsebutton 
%
% Use \selectcolor for a button which opens a color selector popup
\item Select a Color: \selectcolor{Outline color}{yellow}
%
\item Slider: \slider{Transparency}{0}{20mm}
%
% Use \inputtext for a labelled field for user input of text 
\item Input Text: \inputtext{Display Name}{lakes.shp}
\end{itemize}
A shadow indicates a clickable GUI component.

\minisec{Text or Keyboard Conventions}
\begin{itemize}
%
%Use for all urls. Otherwise, it is not clickable in the document.
\item Hyperlinks: \url{http://qgis.org}
%
\item Single Keystroke: press \keystroke{p}
\item Keystroke Combinations: press \keystroke{Ctrl-B}, meaning press and
hold the Ctrl key and then press the B key.
\item Name of a File: \filename{lakes.shp}
%\item Name of a Field: \fieldname{NAMES}
\item Name of a Class: \classname{NewLayer}
\item Method: \method{classFactory}
\item Server: \server{example needed here}
\item SQL Table: \sqltable{example needed here}    
%
%Use usertext for all other text that the user must enter from the keyboard
%that is not covered by any of the above cases
\item User Text: \usertext{qgis ---help}
\end{itemize}

Code is indicated by a fixed-width font:
\begin{verbatim}
PROJCS["NAD_1927_Albers",
  GEOGCS["GCS_North_American_1927",
\end{verbatim}

Platform-specific instructions are indicated as follow.
GUI sequences and small amounts of text can be formatted inline: %
Click \{\nix{}\win{File} \osx{QGIS}\} > Quit to % 
close QGIS. This indicates that on Linux, Unix and Windows platforms, %
click the File menu option first, then Quit from the dropdown menu, while on %
Macintosh OSX platforms, click the QGIS menu option first, then Quit from the %
dropdown menu.

Larger amounts of text may be formatted as a list:
%A moderate amount of text may be formatted as a list:
\begin{itemize}
\item \nix{do this;} 
\item \win{do that;} 
\item \osx{do something else.}
\end{itemize}
or as paragraphs.
%Larger amounts of text should be formatted as paragraphs:

\nix{}Do this and this and this. Then do this
and this and this
and this and this
and this and this
and this and this.

\win{}Do that. Then do that
and that and that
and that and that
and that and that
and that and that
and that and that
and that and that
and that and that.

\osx{}Do something else. 
Then do something else and 
something else and
something else and
something else and
something else and
something else and
something else and
something else and
something else and
something else and
something else altogether.
%Note: on OSX the following commands differ from *NIX and Windows
%- options dialog is called preferences and is found under
% QGIS->Preferences and the shortcut will always be Command , (command
comma)
%- What is usually the Help->about dialog will always be QGIS->About qgis
%- File->Quit will always be QGIS->Quit qgis and the keyboard
shortcut should always be Command Q
%- Colour, font and file dialogs will have differences between platforms
%- if other differences are discovered, please add here




